\chapter{2002年新课标卷（文）}
\section{单选题}
\begin{problem}
直线 \((1 + a)x + y + 1 = 0\) 与圆 \(x^{2} + y^{2} - 2x = 0\) 相切，则 \(a\) 的值为（\quad）

\choices
    {\(\pm 1\)}
    {\(\pm 2\)}
    {\(1\)}
    {\(-1\)}
\end{problem}

\begin{problem}
已知 \(m\)，\(n\) 为异面直线，\(m \subset\) 平面 \(\alpha\)，\(n \subset\) 平面 \(\beta\)，\(\alpha \cap \beta = l\)，则 \(l\)（\quad）

\choices
    {与 \(m, n\) 都相交}
    {与 \(m\)，\(n\) 中至少一条相交}
    {与 \(m\)，\(n\) 都不相交}
    {至多与 \(m\)，\(n\) 中的一条相交}
\end{problem}

\begin{problem}
函数 \( y = a^x \) 在 \([0,1]\) 上的最大值与最小值之和为 3，则 \( a = \)（\quad）

\choices
    {\(\frac{1}{2}\)}
    {\(2\)}
    {\(4\)}
    {\(\frac{1}{4}\)}
\end{problem}

\begin{problem}
不等式 \((1 + x)(1 - |x|) > 0\) 的解集是（\quad）

\choices
    {\(\{x\mid 0\leqslant x < 1\}\)}
    {\(\{x\mid x < 0\) 且 \(x\neq -1\}\)}
    {\(\{x\mid -1 < x < 1\}\)}
    {\(\{x\mid x < 1\) 且 \(x\neq -1\}\)}
\end{problem}

\begin{problem}
在 \((0,2\pi)\) 内，使 \(\sin x > \cos x\) 成立的 \(x\) 的取值范围是（\quad）

\choices
    {\(\left(\frac{\pi}{4},\frac{\pi}{2}\right)\cup \left(\pi ,\frac{5\pi}{4}\right)\)}
    {\(\left(\frac{\pi}{4},\pi\right)\)}
    {\(\left(\frac{\pi}{4},\frac{5\pi}{4}\right)\)}
    {\(\left(\frac{\pi}{4},\pi\right)\cup \left(\frac{5\pi}{4},\frac{3\pi}{2}\right)\)}
\end{problem}

\begin{problem}
设集合 \(M = \left\{x \mid x = \frac{k}{2} + \frac{1}{4}, k \in \Z\right\}\)，\(N = \left\{x \mid x = \frac{k}{4} + \frac{1}{2}, k \in \Z\right\}\)，则（\quad）

\choices
    {\(M = N\)}
    {\(M \subseteq N\)}
    {\(M \supseteq N\)}
    {\(M \cap N = \varnothing\)}
\end{problem}

\begin{problem}
椭圆 \(5x^{2} + ky^{2} = 5\) 的一个焦点是 \((0,2)\)，那么 \(k =\)（\quad）

\choices
    {\(-1\)}
    {\(1\)}
    {\(\sqrt{5}\)}
    {\(-\sqrt{5}\)}
\end{problem}

\begin{problem}
正六棱柱 \(ABCDEF - A_{1}B_{1}C_{1}D_{1}E_{1}F_{1}\) 的底面边长为 1，侧棱长为 \(\sqrt{2}\)，则这个棱柱侧面对角线 \(E_{1}D\) 与 \(BC_{1}\) 所成的角是（\quad）

\choices
    {\(90^{\circ}\)}
    {\(60^{\circ}\)}
    {\(45^{\circ}\)}
    {\(30^{\circ}\)}
\end{problem}

\begin{problem}
函数 \( y = x^{2} + bx + c, x \in [0, +\infty) \) 是单调函数的充要条件是（\quad）

\choices
    {\(b \geqslant 0\)}
    {\(b \leqslant 0\)}
    {\(b > 0\)}
    {\(b < 0\)}
\end{problem}

\begin{problem}
已知 \(0 < x < y < a < 1\)，则有（\quad）

\choices
    {\(\log_a(xy) < 0\)}
    {\(0 < \log_a(xy) < 1\)}
    {\(1 < \log_a(xy) < 2\)}
    {\(\log_a(xy) > 2\)}
\end{problem}

\begin{problem}
从正方体的6个面中选取3个面，其中有2个面不相邻的选法共有（\quad）

\choices
    {\(8\) 种}
    {\(12\) 种}
    {\(16\) 种}
    {\(20\) 种}
\end{problem}

\begin{problem}
平面直角坐标系中，\(O\) 为坐标原点. 已知两点 \(A(3,1)\)，\(B(-1,3)\). 若点 \(C\) 满足 \(\overrightarrow{OC} = \alpha \overrightarrow{OA} + \beta \overrightarrow{OB}\)，其中 \(\alpha\)，\(\beta \in \R\)，且 \(\alpha + \beta = 1\)，则点 \(C\) 的轨迹方程为（\quad）

\choices
    {\( 3x + 2y - 11 = 0 \)}
    {\((x - 1)^{2} + (y - 2)^{2} = 5\)}
    {\(2x - y = 0\)}
    {\(x + 2y - 5 = 0\)}
\end{problem}

\section{填空题}
\begin{problem}
据新华社 2002 年 3 月 12 日电，1985 年到 2000 年间，我国农村人均居住面积如图所示，其中，从\fillinblank{}年\fillinblank{}年的五年间增长最快.

\bitmapfigure[max width=0.58\linewidth,max totalheight=4.8cm]{img/2002/new_curriculum_liberal/q13_fig1.png}
\end{problem}

\begin{problem}
已知 \(\sin 2\alpha = -\sin \alpha\)，\(\alpha \in \left(\frac{\pi}{2}, \pi\right)\)，则 \(\cot \alpha =\)\fillinblank{}.
\end{problem}

\begin{problem}
甲、乙两种冬小麦试验品种连续 5 年的平均单位面积产量如下（单位：\(\mathrm{t}/\mathrm{km}^2\)）：

\begin{center}
\begin{tabular}{|c|c|c|c|c|c|}
\hline
品种 & 第1年 & 第2年 & 第3年 & 第4年 & 第5年 \\
\hline
甲 & \(9.8\) & \(9.9\) & \(10.1\) & \(10\) & \(10.2\) \\
\hline
乙 & \(9.4\) & \(10.3\) & \(10.8\) & \(9.7\) & \(9.8\) \\
\hline
\end{tabular}
\end{center}

其中产量比较稳定的小麦品种是\fillinblank{}.
\end{problem}

\begin{problem}
设函数 \(f(x)\) 在 \((-\infty, +\infty)\) 内有定义，下列函数

\begin{circlelist}
\item \(y = -|f(x)|\)；
\item \(y = x f(x^{2})\)；
\item \(y = -f(-x)\)；
\item \(y = f(x) - f(-x)\).
\end{circlelist}
其中必为奇函数的有\fillinblank{}.（要求填写正确答案的序号）
\end{problem}

\section{解答题}
\begin{problem}
在等比数列 \(\{a_{n}\}\) 中，已知 \(a_6 - a_4 = 24\)，\(a_3a_5 = 64\)，求 \(\{a_{n}\}\) 前8项的和 \(S_{8}\).
\end{problem}

\begin{problem}
已知 \(\sin^2 2\alpha + \sin 2\alpha \cos \alpha - \cos 2\alpha = 1\)，\(\alpha \in \left(0, \frac{\pi}{2}\right)\)，求 \(\sin \alpha\)，\(\tan \alpha\) 的值.
\end{problem}

\begin{problem}
如图，正三棱柱 \(ABC - A_{1}B_{1}C_{1}\) 的底面边长为 \(a\)，侧棱长为 \(\sqrt{2}a\).

\begin{enumerate}
\item 建立适当的坐标系，并写出点 \(A\)，\(B\)，\(A_{1}\)，\(C_{1}\) 的坐标；
\item 求 \(AC_{1}\) 与侧面 \(ABB_{1}A_{1}\) 所成的角.
\end{enumerate}

\bitmapfigure[max width=0.36\linewidth,max totalheight=4.8cm]{img/2002/new_curriculum_liberal/q19_fig1.png}
\end{problem}

\begin{problem}
如图，正方形 \(ABCD\)，\(ABEF\) 的边长都是 \(1\)，而且平面 \(ABCD\)，\(ABEF\) 互相垂直. 点 \(M\) 在 \(AC\) 上移动，点 \(N\) 在 \(BF\) 上移动，若 \(CM = BN = a\) \((0 < a < \sqrt{2})\).

\begin{enumerate}
\item 求 \(MN\) 的长；
\item \(a\) 为何值时，\(MN\) 的长最小；
\item 当 \(MN\) 的长最小时，求面 \(MNA\) 与面 \(MNB\) 所成二面角 \(\alpha\) 的大小.
\end{enumerate}

\bitmapfigure[max width=0.36\linewidth,max totalheight=4.8cm]{img/2002/new_curriculum_liberal/q20_fig1.png}

\end{problem}

\begin{problem}
某单位 6 个员工借助互联网开展工作，每个员工上网的概率都是 \(0.5\)（相互独立）.

\begin{enumerate}
\item 求至少 3 人同时上网的概率；
\item 至少几人同时上网的概率小于 \(0.3\)？
\end{enumerate}
\end{problem}

\begin{problem}
已知 \(a > 0\)，函数 \(f(x) = x^3 - a\)，\(x \in (0, +\infty)\). 设 \(x_1 > 0\)，记曲线 \(y = f(x)\) 在点 \((x_1, f(x_1))\) 处的切线为 \(l\).

\begin{enumerate}
\item 求 \(l\) 的方程；
\item 设 \(l\) 与 \(x\) 轴交点为 \((x_{2}, 0)\). 证明：
\begin{enumerate}
\item \(x_{2} \geqslant a^{\frac{1}{3}}\)；
\item 若 \(x_{2} > a^{\frac{1}{3}}\)，则 \(a^{\frac{1}{3}} < x_{2} < x_{1}\).
\end{enumerate}
\end{enumerate}
\end{problem}

\begin{problem}
已知两点 \(M(-1,0)\)，\(N(1,0)\)，且点 \(P\) 使 \(\overrightarrow{MP} \cdot \overrightarrow{MN}\)，\(\overrightarrow{PM} \cdot \overrightarrow{PN}\)，\(\overrightarrow{NM} \cdot \overrightarrow{NP}\) 成公差小于零的等差数列.

\begin{enumerate}
\item 点 \(P\) 的轨迹是什么曲线？
\item 若点 \(P\) 坐标为 \((x_0, y_0)\)，记 \(\theta\) 为 \(\overrightarrow{PM}\) 与 \(\overrightarrow{PN}\) 的夹角，求 \(\tan \theta\).
\end{enumerate}
\end{problem}
